\chapter*{Заключение}
\addcontentsline{toc}{chapter}{Заключение}

В ходе выполнения первого этапа выпускной квалификационной работы бакалавра были успешно решены все поставленные задачи и получены следующие основные результаты:

\begin{enumerate}
    \item Проведен подробный анализ физико-математических основ одномерной акустики в волноводах переменного сечения на основе уравнения рупора Вебстера и граничных условий.
    \item Разработан и программно реализован в виде единого Python/C++ комплекса набор численных решателей, включающий:
    \begin{itemize}
        \item метод матриц передачи (TMM) для кусочно-цилиндрической и кусочно-конической аппроксимаций геометрии волновода;
        \item авторегрессионный z-доменный метод моделирования (ARMA);
        \item метод матриц длинных линий (TLM) во временной области в цилиндрической формулировке на C++.
    \end{itemize}
    \item На тестовых задачах регулярных форм (цилиндр, конус) проведена верификация численных схем, показавшая высокую точность расчетов и полное физическое соответствие.
    \item Исследовано качество геометрической аппроксимации волновода. Установлено, что коническая дискретизация обеспечивает меньшую относительную ошибку площади сечения по сравнению со ступенчатой цилиндрической при одинаковом числе разбиений.
    \item Выполнен сравнительный анализ точности расчета передаточных функций. Выявлено преимущество конического TMM-решателя для гладких геометрий, а также высокая точность быстрого авторегрессионного решателя (ARMA).
    \item Выполнена сравнительная оценка вычислительной производительности реализованных алгоритмов. Частотные методы (TMM, ARMA) показали наилучшую скорость расчетов на индивидуальных геометриях по сравнению с временным методом TLM.
\end{enumerate}

Полученные результаты и разработанный программный комплекс численных решателей могут служить основой для быстрого моделирования и массовой генерации данных в задачах одномерного акустического анализа. В качестве перспектив развития работы планируется интеграция реализованных физических решателей с нейросетевыми операторными моделями (FNO) для ускорения решения прямой задачи и оптимизационными методами для решения обратной задачи восстановления геометрии.
